\chapter{Gilbert Syndrome}

\section*{Definition and Overview}
Gilbert syndrome is a common, harmless genetic condition in which the liver processes bilirubin slightly less efficiently than usual, leading to mildly elevated levels of bilirubin in the blood. Bilirubin is a yellow pigment produced when red blood cells are broken down, and it is normally removed by the liver. In Gilbert syndrome, a reduced amount of the enzyme that conjugates bilirubin means that levels can rise, especially during stress, fasting, illness, or exertion. This causes mild jaundice (a slight yellowing of the skin and eyes) in some people. Gilbert syndrome is not a disease, causes no liver damage, and requires no treatment; its main importance is that it is often discovered by accident and should be distinguished from more serious causes of jaundice.

\section*{Epidemiology}
Gilbert syndrome is one of the most common inherited conditions, affecting around 5--10\% of the population, although many people never know they have it. It is more common in males than females, with estimates suggesting men are two to three times more likely to be affected. The condition occurs across all ethnic groups, though prevalence varies. It is usually diagnosed in young adulthood, often when routine blood tests, taken for unrelated reasons, show a mildly elevated bilirubin level. Because it is harmless, its true prevalence is probably higher than diagnosed.

\section*{Aetiology and Pathophysiology}
Gilbert syndrome is caused by a genetic variation that reduces the activity of the enzyme responsible for processing bilirubin.

\begin{itemize}
    \item \textbf{Genetic cause:} a variation in the \textit{UGT1A1} gene, inherited in an autosomal recessive pattern, reduces the activity of the enzyme UDP-glucuronosyltransferase
    \item \textbf{Reduced conjugation:} the liver conjugates bilirubin less efficiently, so unconjugated (indirect) bilirubin accumulates in the blood
    \item \textbf{Mild elevation:} bilirubin levels are usually only mildly raised, typically below 50 micromol/L, and can fluctuate
    \item \textbf{Triggers:} fasting, dehydration, physical exertion, stress, illness, infection, and alcohol can temporarily raise bilirubin levels further
    \item \textbf{No liver damage:} the liver itself is structurally and functionally normal, and there is no inflammation or fibrosis
    \item \textbf{Benign course:} the condition has no effect on life expectancy or liver health
\end{itemize}

\section*{Clinical Features}
Most people with Gilbert syndrome have no symptoms at all and are diagnosed by chance.

\begin{itemize}
    \item A mildly yellow tinge to the skin or whites of the eyes, usually only noticeable during illness, fasting, or stress
    \item Fatigue, which some people report, although it is not clearly caused by the condition
    \item Elevated bilirubin on blood tests performed for other reasons
    \item No pain, itching, or other symptoms of liver disease
    \item Symptoms, if any, are intermittent and generally mild
    \item Normal energy levels and general health between episodes
\end{itemize}

\section*{Differential Diagnosis}
Because jaundice can indicate serious disease, Gilbert syndrome must be distinguished from other causes of elevated bilirubin.

\begin{itemize}
    \item Haemolysis: increased breakdown of red blood cells, which also raises unconjugated bilirubin and requires investigation
    \item Liver disease, including viral hepatitis, alcoholic liver disease, and cirrhosis
    \item Other inherited conditions of bilirubin metabolism, such as Crigler-Najjar syndrome, which is far more severe
    \item Biliary obstruction, such as gallstones or tumours, which raises conjugated bilirubin and is associated with dark urine and pale stools
    \item Drug-induced liver injury
    \item Thyrotoxicosis and other conditions that can mildly elevate bilirubin
    \item Fasting-related jaundice in an otherwise healthy person
\end{itemize}

\section*{When to Seek Medical Care}
Gilbert syndrome itself requires no treatment, but anyone with yellowing of the skin or eyes, especially if new, persistent, or accompanied by other symptoms, should see a doctor to confirm the cause. People who know they have Gilbert syndrome should mention it when having blood tests so results are interpreted correctly.

\textbf{Call 000 (or local emergency services) for:}
\begin{itemize}
    \item Jaundice with severe abdominal pain, high fever, or confusion, which may indicate a serious condition such as biliary infection or liver failure
    \item Dark urine and pale stools with jaundice, which suggest obstruction and need urgent review
    \item Vomiting blood or signs of internal bleeding
    \item Sudden severe illness in anyone with known liver disease
\end{itemize}

\section*{Diagnosis and Investigations}
Gilbert syndrome is diagnosed by a combination of typical blood test findings and exclusion of other causes.

\begin{itemize}
    \item \textbf{Blood tests:} mildly elevated unconjugated bilirubin with normal liver enzymes (ALT, AST, ALP), normal albumin, and normal full blood count
    \item \textbf{Reticulocyte count and haptoglobin:} exclude haemolysis as a cause of the raised bilirubin
    \item \textbf{Response to triggers:} bilirubin may rise with fasting and fall when eating normally, supporting the diagnosis
    \item \textbf{Genetic testing:} rarely needed in routine practice, but can confirm the \textit{UGT1A1} variation if the diagnosis is uncertain
    \item \textbf{Exclusion of liver disease:} ultrasound and further tests are arranged only if liver function is otherwise abnormal or symptoms suggest another cause
\end{itemize}

\section*{Management}
There is no treatment for Gilbert syndrome because it is not a disease and causes no harm.

\begin{itemize}
    \item \textbf{Reassurance:} the mainstay of management is explaining that the condition is harmless and requires no treatment
    \item \textbf{Avoiding unnecessary tests:} a confirmed diagnosis prevents repeated, needless investigations for the mild jaundice
    \item \textbf{Avoiding triggers:} regular meals, adequate hydration, and avoiding excessive fasting can reduce episodes of noticeable jaundice
    \item \textbf{Medicine awareness:} a small number of medicines are metabolised by the same enzyme, and their effects can differ slightly in Gilbert syndrome; this matters only for a few drugs, and doctors can advise
    \item \textbf{Normal life:} no restrictions on diet, activity, travel, or employment are needed
\end{itemize}

\section*{Complications}
\begin{itemize}
    \item There are no medical complications of Gilbert syndrome itself
    \item The main "complications" are psychosocial: anxiety from a confusing diagnosis, and unnecessary investigations if the condition is not recognised
    \item Rarely, people are denied life insurance or asked about their liver condition, though the condition is benign
    \item Slightly increased risk of gallstones in some studies, though the practical significance is small
\end{itemize}

\section*{Prognosis and Outcomes}
Gilbert syndrome has an excellent prognosis. It does not affect life expectancy, does not progress to liver disease, and requires no treatment. Bilirubin levels may fluctuate throughout life but never reach levels that cause harm. The most important outcome of diagnosis is peace of mind: once the condition is properly identified, people can be reassured that their "abnormal" blood test is normal for them and requires no further concern or follow-up.

\section*{Prevention and Self-Care}
\begin{itemize}
    \item If you have Gilbert syndrome, carry a note or mention it when blood tests are ordered
    \item Eat regular meals and avoid prolonged fasting, which can raise bilirubin
    \item Keep well hydrated, especially during illness or exertion
    \item Understand that mild jaundice during stress or illness is expected and harmless
    \item See a doctor if jaundice persists, becomes severe, or is accompanied by other symptoms, to rule out other causes
    \item Do not restrict your diet, activity, or lifestyle because of the diagnosis
    \item Reassure family members, since the condition can run in families
\end{itemize}

\section*{Sources and Further Reading}
The following reputable sources provide further information for healthcare students and patients seeking reliable, current advice about Gilbert syndrome.

\begin{itemize}
    \item Healthdirect Australia. \textit{Gilbert syndrome.} https://www.healthdirect.gov.au/
    \item Australian Liver Foundation. \textit{Information about Gilbert syndrome.} https://www.liver.org.au/
    \item NHS. \textit{Gilbert's syndrome.} https://www.nhs.uk/
    \item Mayo Clinic. \textit{Gilbert syndrome: symptoms and causes.} https://www.mayoclinic.org/
    \item MedlinePlus. \textit{Gilbert syndrome.} https://medlineplus.gov/
    \item MSD Manual. \textit{Gilbert syndrome.} https://www.msdmanuals.com/
\end{itemize}
